\chapter{Những Điều Về Thất Bại}
\markboth{Những Điều Về Thất Bại}{Things a Teacher Wants to Say About Failure}

\section{Ai Cũng Từng Thất Bại}
\begin{truyen}{Ai Cũng Từng Thất Bại}{Everyone Has Failed at Something}
\chuhoa{N}{ếu em nhìn đủ kỹ, em sẽ thấy hầu như ai cũng từng thất bại ở một chỗ nào đó.}
\textit{(If you look carefully enough, you will see that almost everyone has failed somewhere.)}

Chỉ là không phải ai cũng kể ra những đoạn đó.
\textit{(The difference is that not everyone tells those stories aloud.)}

Vì vậy học sinh rất dễ tưởng chỉ mình mình đang vấp.
\textit{(So students easily imagine they are the only ones stumbling.)}

Biết rằng thất bại là chuyện con người nào cũng đi qua giúp em bớt xấu hổ hơn với đoạn đường của mình.
\textit{(Knowing failure belongs to every human path can reduce shame about your own.)}
\end{truyen}

\begin{baihoc}
Thất bại không làm em thành trường hợp dị thường. Nó chỉ cho thấy em cũng đang đi trên đường lớn lên như mọi người.
\textit{(Failure does not make you abnormal. It shows that you are also walking the path of growth.)}
\end{baihoc}

\begin{ghinhoanh}
\textbf{Short English to remember:}

Failure is common in every human journey.

\end{ghinhoanh}

\begin{tuhocnhanh}
1. Vì sao học sinh hay tưởng chỉ mình mình thất bại? \textit{Why do students often think only they are failing?}

2. Biết ai cũng từng vấp giúp em ở đâu? \textit{How does it help to know everyone stumbles?}

3. Em đang xấu hổ quá mức về điều gì? \textit{What are you feeling too ashamed about?}
\end{tuhocnhanh}
\ngancach

\section{Sai Không Đáng Xấu Hổ}
\begin{truyen}{Sai Không Đáng Xấu Hổ}{Being Wrong Is Not Shameful}
\chuhoa{C}{ái đáng xấu hổ không phải lúc nào cũng là sai. Có khi cái đáng ngại hơn là giả vờ đúng khi mình chưa hiểu.}
\textit{(What is shameful is not always being wrong. Sometimes the more troubling thing is pretending to be right when you do not understand.)}

Sai là tín hiệu cho biết còn chỗ phải học thêm.
\textit{(Being wrong is a signal that more learning is needed.)}

Nếu em đủ bình tĩnh để nhìn vào nó, cái sai trở thành bản đồ chỉ ra nơi mình cần sửa.
\textit{(If you are calm enough to look at it, a mistake becomes a map showing where repair is needed.)}

Xấu hổ quá mức với cái sai chỉ làm em muốn trốn khỏi bài học của chính mình.
\textit{(Excessive shame about mistakes only makes you want to hide from your own lesson.)}
\end{truyen}

\begin{baihoc}
Sai là một phần của học thật. Nó không đáng xấu hổ bằng việc sống mãi trong sự giả vờ.
\textit{(Mistakes belong to real learning. They are less shameful than living in pretense.)}
\end{baihoc}

\begin{ghinhoanh}
\textbf{Short English to remember:}

A mistake is easier to fix than a false image.

\end{ghinhoanh}

\begin{tuhocnhanh}
1. Vì sao cái sai không nên bị coi là nhục? \textit{Why should mistakes not be treated as humiliation?}

2. Điều gì còn hại hơn sai? \textit{What can be more harmful than being wrong?}

3. Em đang trốn điều gì vì sợ sai? \textit{What are you avoiding because of fear of mistakes?}
\end{tuhocnhanh}
\ngancach

\section{Thất Bại Giúp Mình Mạnh Hơn}
\begin{truyen}{Thất Bại Giúp Mình Mạnh Hơn}{Failure Can Make Us Stronger}
\chuhoa{T}{ất nhiên thất bại tự nó không thần kỳ. Nó chỉ làm mình mạnh hơn khi mình chịu học từ nó.}
\textit{(Of course failure is not magical by itself. It only makes us stronger when we learn from it.)}

Nhưng đúng là sau những lần vấp, con người có thể bớt ảo tưởng hơn, bền hơn và hiểu mình thật hơn.
\textit{(But it is true that after setbacks, people can become less illusioned, more resilient, and more honest about themselves.)}

Nhiều sức mạnh sâu không đến từ những ngày suôn sẻ, mà đến từ lúc mình bị ép phải lớn lên.
\textit{(Many deeper strengths come not from smooth days, but from seasons that force us to grow.)}

Vấn đề là đừng chỉ đau rồi thôi. Hãy đau xong rồi nhìn lại.
\textit{(The key is not to hurt and stop there. Hurt, then reflect.)}
\end{truyen}

\begin{baihoc}
Thất bại không tự động làm em mạnh. Nhưng nếu em nhìn lại đúng cách, nó có thể làm em trưởng thành hơn thật.
\textit{(Failure does not automatically strengthen you. But if you reflect properly, it can genuinely mature you.)}
\end{baihoc}

\begin{ghinhoanh}
\textbf{Short English to remember:}

Pain becomes strength only when it is understood.

\end{ghinhoanh}

\begin{tuhocnhanh}
1. Vì sao không phải thất bại nào cũng tự làm con người tốt hơn? \textit{Why does failure not automatically improve a person?}

2. Điều gì cần xảy ra sau một lần vấp để mình mạnh hơn? \textit{What needs to happen after failure for strength to grow?}

3. Em đã học được gì từ một lần thất bại gần đây? \textit{What have you learned from a recent failure?}
\end{tuhocnhanh}
\ngancach

\section{Đừng Bỏ Cuộc Sau Một Lần Sai}
\begin{truyen}{Đừng Bỏ Cuộc Sau Một Lần Sai}{Do Not Quit After One Mistake}
\chuhoa{M}{ột lần sai không đủ để kết luận rằng em không hợp, không làm nổi, hay không có tương lai ở việc đó.}
\textit{(One mistake is not enough to conclude that you are not suited, incapable, or without a future in something.)}

Nhưng rất nhiều học sinh dừng lại quá sớm.
\textit{(Yet many students stop too early.)}

Các em biến một lần vấp thành điểm kết thúc.
\textit{(They turn one stumble into the end point.)}

Trong khi đó, nhiều điều đáng giá chỉ mở ra cho người chịu thử lại thêm vài lần nữa.
\textit{(Meanwhile, many worthwhile things open only to those who try again a few more times.)}
\end{truyen}

\begin{baihoc}
Đừng để một lần sai lấy mất khỏi em cơ hội thấy phiên bản tốt hơn của mình ở lần sau.
\textit{(Do not let one mistake steal your chance to see a better version of yourself next time.)}
\end{baihoc}

\begin{ghinhoanh}
\textbf{Short English to remember:}

One mistake should not become the end of the story.

\end{ghinhoanh}

\begin{tuhocnhanh}
1. Vì sao học sinh hay bỏ sau một lần sai? \textit{Why do students often quit after one mistake?}

2. Điều gì có thể mở ra nếu em thử thêm vài lần nữa? \textit{What may open if you try a few more times?}

3. Em đang muốn bỏ điều gì quá sớm? \textit{What are you quitting too early?}
\end{tuhocnhanh}
\ngancach

\section{Người Thành Công Từng Thất Bại Nhiều Lần}
\begin{truyen}{Người Thành Công Từng Thất Bại Nhiều Lần}{Successful People Often Failed Many Times}
\chuhoa{N}{gười ta thường chỉ nhìn phần thành công đã hiện ra mà quên mất đoạn trước đó rất lộn xộn.}
\textit{(People usually see the visible success and forget the messy road before it.)}

Nhiều người thành công đi qua rất nhiều lần không đạt, không được chọn, không được công nhận.
\textit{(Many successful people went through repeated failures, rejection, and lack of recognition.)}

Điều khác biệt không chỉ nằm ở tài năng, mà còn ở việc họ không để mỗi lần vấp chặn luôn đường đi.
\textit{(The difference is not only talent, but also refusing to let each setback block the whole path.)}

Biết điều này không phải để em lạc quan mù quáng, mà để em đỡ tuyệt vọng quá sớm.
\textit{(Knowing this is not for blind optimism, but to keep you from giving up too early.)}
\end{truyen}

\begin{baihoc}
Con đường đi tới điều tốt hiếm khi thẳng. Vấp nhiều lần không có nghĩa là đường đã hết.
\textit{(The road to something good is rarely straight. Multiple setbacks do not mean the road is over.)}
\end{baihoc}

\begin{ghinhoanh}
\textbf{Short English to remember:}

Repeated failure does not prove the road is closed.

\end{ghinhoanh}

\begin{tuhocnhanh}
1. Vì sao ta hay quên đoạn thất bại của người thành công? \textit{Why do we forget the failures behind success?}

2. Điều gì giúp họ đi tiếp? \textit{What helps them continue?}

3. Em đang tuyệt vọng quá sớm ở chỗ nào? \textit{Where are you becoming hopeless too early?}
\end{tuhocnhanh}
\ngancach

\section{Thất Bại Dạy Điều Sách Không Dạy}
\begin{truyen}{Thất Bại Dạy Điều Sách Không Dạy}{Failure Teaches What Books Cannot}
\chuhoa{S}{ách dạy kiến thức, nhưng có những điều chỉ khi vấp con người mới thật sự hiểu.}
\textit{(Books teach knowledge, but some things are understood only after stumbling.)}

Thất bại dạy khiêm tốn, dạy cách chịu trách nhiệm, dạy mình đã yếu ở đâu và dạy cả cách đứng dậy thế nào.
\textit{(Failure teaches humility, responsibility, where we are weak, and how to rise again.)}

Không cuốn sách nào truyền những bài đó vào mình sâu bằng một lần mình phải tự đi qua nó.
\textit{(No book can plant those lessons as deeply as living through them.)}

Cho nên đừng chỉ nhìn thất bại như mất mát. Nhiều khi nó còn là một lớp học riêng của đời.
\textit{(So do not see failure only as loss. Sometimes it is one of life's private classrooms.)}
\end{truyen}

\begin{baihoc}
Có những bài học chỉ thành thật khi đi qua va đập, chứ không chỉ qua lời giảng.
\textit{(Some lessons become real only through collision, not through explanation alone.)}
\end{baihoc}

\begin{ghinhoanh}
\textbf{Short English to remember:}

Some truths are learned only through being tested.

\end{ghinhoanh}

\begin{tuhocnhanh}
1. Vì sao thất bại dạy được điều sách không dạy? \textit{Why can failure teach what books cannot?}

2. Những bài học nào em thường chỉ hiểu sau khi vấp? \textit{What lessons do people often understand only after failing?}

3. Em có đang nhìn thất bại chỉ như mất mát không? \textit{Are you seeing failure only as loss?}
\end{tuhocnhanh}
\ngancach

\section{Học Từ Lỗi Sai}
\begin{truyen}{Học Từ Lỗi Sai}{Learn from Your Errors}
\chuhoa{L}{ỗi sai chỉ thật sự có giá trị khi em chịu nhìn kỹ nó.}
\textit{(Errors become valuable only when you are willing to examine them carefully.)}

Làm sai xong bỏ qua thì cái sai ấy rất dễ quay lại.
\textit{(If you make a mistake and move on carelessly, it often returns.)}

Nhưng nếu em tìm ra vì sao sai, sai ở bước nào, do hổng kiến thức hay do vội, em sẽ học được nhiều hơn từ chính lỗi đó.
\textit{(But if you find why it happened, at which step, and whether it came from a gap or from rushing, you learn far more from it.)}

Sửa lỗi là một kỹ năng. Và người giỏi thường giỏi kỹ năng này.
\textit{(Correcting errors is a skill, and strong learners are often strong at it.)}
\end{truyen}

\begin{baihoc}
Đừng chỉ buồn vì lỗi sai. Hãy lấy nó ra làm bài học cụ thể cho lần sau.
\textit{(Do not only feel sad about mistakes. Turn them into a concrete lesson for next time.)}
\end{baihoc}

\begin{ghinhoanh}
\textbf{Short English to remember:}

A mistake becomes useful when it is studied.

\end{ghinhoanh}

\begin{tuhocnhanh}
1. Vì sao lỗi sai dễ lặp lại nếu không nhìn kỹ? \textit{Why do mistakes repeat when we do not examine them?}

2. Em nên hỏi gì sau khi làm sai? \textit{What should you ask after making a mistake?}

3. Một lỗi cũ nào em cần ngồi lại nhìn cho rõ? \textit{What old error do you need to study carefully?}
\end{tuhocnhanh}
\ngancach

\section{Sai Hôm Nay, Giỏi Ngày Mai}
\begin{truyen}{Sai Hôm Nay, Giỏi Ngày Mai}{Wrong Today, Better Tomorrow}
\chuhoa{M}{ột người học giỏi không phải lúc nào cũng đã giỏi sẵn. Nhiều khi họ chỉ là người đi qua nhiều lần chưa giỏi mà không bỏ cuộc.}
\textit{(A strong learner is not always someone who started strong. Often they are someone who survived many not-yet-good days without quitting.)}

Sai hôm nay không khóa luôn khả năng giỏi ngày mai.
\textit{(Being wrong today does not lock away the chance of being better tomorrow.)}

Miễn là em còn học, còn sửa, còn luyện, ngày mai vẫn có thể khác.
\textit{(As long as you keep learning, correcting, and practicing, tomorrow can still be different.)}

Điều quan trọng là đừng dùng hôm nay để kết án ngày mai quá sớm.
\textit{(What matters is not using today to condemn tomorrow too early.)}
\end{truyen}

\begin{baihoc}
Ngày mai tốt hơn thường được xây trên cách em đối xử với cái sai của hôm nay.
\textit{(A better tomorrow is often built on how you handle today's mistakes.)}
\end{baihoc}

\begin{ghinhoanh}
\textbf{Short English to remember:}

Today\'s mistake does not cancel tomorrow\'s growth.

\end{ghinhoanh}

\begin{tuhocnhanh}
1. Vì sao hôm nay chưa tốt không quyết định luôn ngày mai? \textit{Why does a weak today not decide tomorrow?}

2. Điều gì giúp ngày mai khác đi? \textit{What helps tomorrow become different?}

3. Em đang phán xét tương lai mình quá sớm ở đâu? \textit{Where are you judging your future too early?}
\end{tuhocnhanh}
\ngancach

\section{Đừng Sợ Thử Lại}
\begin{truyen}{Đừng Sợ Thử Lại}{Do Not Be Afraid to Try Again}
\chuhoa{T}{hử lại là một trong những việc khó nhất sau khi em đã từng đau vì thất bại.}
\textit{(Trying again is one of the hardest things after being hurt by failure.)}

Lần trước càng đau, lần này càng dễ run.
\textit{(The more painful the last time was, the more you tremble this time.)}

Nhưng chính hành động thử lại mới cho em cơ hội sửa điều cũ và tạo ra một kết quả khác.
\textit{(But trying again is exactly what gives you a chance to correct the past and create a different result.)}

Thầy không nói thử lại lúc nào cũng thắng. Thầy chỉ nói không thử lại thì chắc chắn không có cửa nào mở thêm.
\textit{(I am not saying retrying always wins. I am saying that without retrying, no new door opens.)}
\end{truyen}

\begin{baihoc}
Thử lại là cách con người từ chối để một lần đau định nghĩa luôn khả năng của mình.
\textit{(Trying again is how people refuse to let one painful failure define their ability forever.)}
\end{baihoc}

\begin{ghinhoanh}
\textbf{Short English to remember:}

Trying again keeps possibility alive.

\end{ghinhoanh}

\begin{tuhocnhanh}
1. Vì sao thử lại rất khó sau một lần đau? \textit{Why is retrying hard after a painful failure?}

2. Không thử lại lấy đi của em điều gì? \textit{What does refusing to retry take away from you?}

3. Em đang cần thử lại điều gì? \textit{What do you need to try again?}
\end{tuhocnhanh}
\ngancach

\section{Thất Bại Là Một Bài Học}
\begin{truyen}{Thất Bại Là Một Bài Học}{Failure Is Also a Lesson}
\chuhoa{S}{au cùng, nếu nhìn đúng, thất bại là một dạng bài học rất đắt giá.}
\textit{(In the end, if seen correctly, failure is a costly but valuable lesson.)}

Nó không vui, nhưng nó thật.
\textit{(It is not pleasant, but it is real.)}

Nó ép em đối diện với giới hạn, với lỗi của mình, với chỗ mình còn non.
\textit{(It forces you to face your limits, your mistakes, and the places where you are still immature.)}

Nếu em học được từ nó, em sẽ không chỉ giỏi hơn. Em còn thật hơn và vững hơn.
\textit{(If you learn from it, you will not only become better. You will also become truer and steadier.)}
\end{truyen}

\begin{baihoc}
Một thất bại được hiểu đúng có thể trở thành một trong những bài học tốt nhất của đời học sinh.
\textit{(A failure understood rightly can become one of the best lessons in a student's life.)}
\end{baihoc}

\begin{ghinhoanh}
\textbf{Short English to remember:}

Failure can become a teacher when you stop running from it.

\end{ghinhoanh}

\begin{tuhocnhanh}
1. Vì sao thất bại cũng là một bài học? \textit{Why is failure also a lesson?}

2. Nó ép con người đối diện với điều gì? \textit{What does it force a person to face?}

3. Em muốn lấy bài học nào ra từ một lần thất bại của mình? \textit{What lesson do you want to extract from one failure?}
\end{tuhocnhanh}
\ngancach
